%Brenon-Audat, pg 241

Le mélange \ce{Cu}-\ce{Ni} cristallise en formant une phase homogène, en apparence homogène, lors du refroidissement théorique. L'observation microscopique d'une composition \ce{Cu}:\ce{Ni} 7:3 montre l'existence d'hétérogèneité au sein d'un grain.

\begin{center}
\begin{minipage}{0.48\textwidth}
\centering
\includegraphics[width=0.5\textwidth]{/home/jeff/Documents/IUTOrsay/1-Cours/4p13-Alliages/Cours-TD/n-3-Lecture&Utilisation_Binaires&Ternaires/Chp3-Ex/Homogene-Ex02/figures/Fig01}
\end{minipage}\hfill\begin{minipage}{0.48\textwidth}
\captionof{figure}{Solution solide Ni/Cu observée au microscope optique, préalablement traité par un mélange \ce{CH4OH}/\ce{H2O2}. La cristalliation a lieu sous forme de dendrite}
\end{minipage}
\end{center}

L'analyse de la composition de l'échantillon montre qu'au coeur des grains de nickel, la fraction massique vaut $0.8$, tandis qu'aux joints, elle approche de $0.2$. Sachant que le système \ce{Cu}/\ce{Ni} présente, à l'équilibre, un diagramme à fuseau idéal, expliquer l'hétérogeneité observée et justifier l'intérêt d'un recuit afin d'homogeneiser le milieu.

\paragraph*{Données}
\begin{align*}
T^\star_{fus} (\ce{Ni}) = 1453 ~ \si{\kelvin} ~ ; ~ T^\star_{fus} (\ce{Cu}) = 1083 ~ \si{\kelvin}
\end{align*}